\section{Prediction intervals, detection, and the sizing problem}
\label{sec:sizing}

\subsection{Calibrated prediction intervals}

A point estimate of throughput at an unmeasured validator count is not
actionable: capacity planning needs a lower bound holding with a stated
probability. Two sources of uncertainty matter --- parameter uncertainty in
$(\gamma,\beta)$, which grows with the extrapolation distance, and residual
dispersion, which does not. With $\mathbf{z}=(1,\log\cc,-\log\n)^{\!\top}$ and
$\hat\vartheta=(\log\Tzero,\gamma,\beta)^{\!\top}$, the delta
method~\cite{seber2003} gives the log-space predictive variance
$\hat{V}(\n,\cc)=\mathbf{z}^{\!\top}\widehat{\Cov}(\hat\vartheta)\mathbf{z}
+\hat\sigma^{2}$.

\begin{State}[Lognormal prediction interval]
\label{st:pi}
If the log-space residuals are approximately normal, an approximate
$(1-\epsp)$ prediction interval for $\TPS(\n,\cc)$ is
$\exp\{\log\widehat{\TPS}(\n,\cc)\mp\zq_{1-\epsp/2}\hat{V}^{1/2}(\n,\cc)\}$,
with one-sided lower endpoint $\underline{\TPS}_{1-\epsp}(\n,\cc)$ obtained by
replacing $\zq_{1-\epsp/2}$ with $\zq_{1-\epsp}$.
\end{State}

Normality of log-residuals is an assumption and heteroscedasticity across $\n$
is expected, so we do not rely on it: a residual bootstrap~\cite{efron1993}
resamples log-residuals, refits~\eqref{eq:ols} and reports empirical quantiles
of the predictive distribution. The interval is then validated rather than
asserted --- fitting on $\n\in\{4,6,8,10\}$ and measuring the fraction of
held-out observations at $\n\in\{12,16\}$ falling inside the nominal $90\%$ and
$95\%$ intervals. Coverage substantially below nominal falsifies the interval,
and we report it whichever way it comes out; this replaces the mean absolute
percentage error customary in this literature, which cannot be checked against a
target.

\subsection{Detection of coordinated omission faults}

The bound below is Hoeffding's inequality with a design-effect correction and we
claim no novelty for it; its role is to supply a constraint whose monotonicity
in~$\n$ opposes that of throughput, with parameters we measure rather than
postulate.

\begin{Def}[Fault class and detector]
\label{def:faults}
A validator subset $B$ with $|B|=\kf$ is \emph{faulty} if each member fails to
contribute to a proposal round independently with probability $\pd>\pf$, where
$\pf$ is the corresponding rate for a correct validator; this is the
\emph{omission} class, realized at the network level and excluding equivocation.
The detector observes each validator's participation over a window of $\w$
rounds and flags~$i$ when its participation rate falls below a threshold
$\thr\in(\pf,\pd)$; detection succeeds when at least one member of $B$ is
flagged.
\end{Def}

For a single validator, Hoeffding's inequality~\cite{hoeffding1963} bounds the
probability that $\w$ independent observations place the empirical rate on the
wrong side of $\thr$ by $\exp\{-2\w(\pd-\thr)^{2}\}$. Validators controlled by
one adversary do not, however, fail independently.

\begin{Assumption}[Exchangeable correlation]
\label{as:rho}
Participation indicators of validators in $B$ are exchangeable with pairwise
correlation $\rho\in[0,1)$, so the variance of their mean carries the design
effect $1+(\kf-1)\rho$.
\end{Assumption}

Exchangeability is a modelling choice, not a fact about adversaries: a real
attacker may correlate failures in a structured, time-varying way that no single
$\rho$ describes. We adopt it because it yields a closed form with one
interpretable parameter, we estimate $\rho$ from injected faults rather than
assuming a value, and we treat the result as a design heuristic calibrated on
the omission class rather than as a worst-case guarantee.

\begin{State}[Miss probability and the minimum validator count]
\label{st:rho}
Under Definition~\ref{def:faults} and Assumption~\ref{as:rho}, with
$\kf=\lceil\phi\n\rceil$ for an assumed faulty fraction $\phi$,
\begin{equation}
\pmiss(\n,\w)\le\exp\!\left\{-\frac{2\w\kf(\pd-\thr)^{2}}{1+(\kf-1)\rho}\right\}.
\label{eq:pmiss}
\end{equation}
The bound is nonincreasing in $\n$ and in $\w$, and $\pmiss(\n,\w)\le\epss$
holds for all $\n\ge\nmin$ with
\begin{equation}
\nmin=\left\lceil\frac{1}{\phi}\cdot
\frac{\bigl(1+(\kf-1)\rho\bigr)\log(1/\epss)}{2\w(\pd-\thr)^{2}}\right\rceil .
\label{eq:nmin}
\end{equation}
\end{State}

\begin{Proof}
A miss requires every member of $B$ to remain unflagged. Under
Assumption~\ref{as:rho} the effective number of independent observations is
$\neff=\kf\w/(1+(\kf-1)\rho)$, and Hoeffding's inequality applied to the
exchangeable mean with this effective sample size gives~\eqref{eq:pmiss}. The
exponent is proportional to $\kf/(1+(\kf-1)\rho)$, whose derivative in $\kf$ has
the sign of $1-\rho>0$; since $\kf$ is nondecreasing in $\n$ and the exponent is
increasing in $\w$, the bound is nonincreasing in both. Solving
$\pmiss\le\epss$ for $\kf$ and substituting $\kf=\phi\n$ yields~\eqref{eq:nmin},
and monotonicity makes the constraint a half-line.
\end{Proof}

\subsection{Chance-constrained sizing and the feasibility window}

\begin{Problem}[Validator-set sizing]
\label{prob:sizing}
Given a peak demand $\Dpk$, a confidence level $1-\epsp$, a safety target
$\epss$, a window $\w$, a faulty fraction $\phi$ and a nondecreasing cost
$\Cost(\n)$, find
\begin{equation}
\nopt\in\arg\min_{\n\in\mathbb{N}}\ \Cost(\n)
\quad\text{s.t.}\quad
\Prob\bigl[\TPS(\n,\cc^{\star})\ge\Dpk\bigr]\ge1-\epsp,
\quad \pmiss(\n,\w)\le\epss ,
\label{eq:sizing}
\end{equation}
where $\cc^{\star}=\arg\max_{\cc}\gfun(\cc)$.
\end{Problem}

\begin{Lemma}[Deterministic equivalent and monotonicity]
\label{lem:detequiv}
The chance constraint in~\eqref{eq:sizing} is equivalent to
\begin{equation}
h(\n)\eqdef\log\Tzero+\log\gfun(\cc^{\star})-\beta\log\n
-\zq_{1-\epsp}\hat{V}^{1/2}(\n,\cc^{\star})-\log\Dpk\ \ge 0 .
\label{eq:detequiv}
\end{equation}
If $\hat{V}$ is nondecreasing in $\log\n$ over the extrapolation region, then
$h$ is strictly decreasing there.
\end{Lemma}

\begin{Proof}
By Statement~\ref{st:pi} the event $\TPS\ge\Dpk$ has probability at least
$1-\epsp$ exactly when $\underline{\TPS}_{1-\epsp}\ge\Dpk$; taking logarithms
gives~\eqref{eq:detequiv}. The first three terms strictly decrease in $\log\n$
because $\beta>0$, and $-\zq\hat{V}^{1/2}$ is nonincreasing by hypothesis. That
hypothesis holds whenever the design is centred at or below the operating range,
since $\hat{V}$ is quadratic in $\log\n$ with minimum at the design centroid.
\end{Proof}

\begin{Theorem}[Feasibility window]
\label{thm:window}
Under the hypotheses of Lemma~\ref{lem:detequiv} and Statement~\ref{st:rho}, the
feasible set of Problem~\ref{prob:sizing} is the integer interval
$\Feas=\{\n:\nmin\le\n\le\nmax\}$ with $\nmin$ given by~\eqref{eq:nmin} and
$\nmax=\lfloor\sup\{\n:h(\n)\ge0\}\rfloor$. If $\Feas\ne\emptyset$ and $\Cost$
is nondecreasing then $\nopt=\nmin$, and both endpoints are computable in
$O(\log\nmax)$ evaluations of $h$ by bisection.
\end{Theorem}

\begin{Proof}
By Statement~\ref{st:rho} the detection constraint holds exactly on
$\{\n\ge\nmin\}$; by Lemma~\ref{lem:detequiv} $h$ is strictly decreasing, so
$\{\n:h(\n)\ge0\}$ is a down-set and the throughput constraint holds exactly on
$\{\n\le\nmax\}$. The feasible set is the intersection of two opposing
half-lines, hence an interval; it is finite, so a minimizer exists whenever it
is nonempty, and monotonicity of $\Cost$ places it at the left endpoint.
Bisection on the monotone $h$ locates $\nmax$ in $O(\log\nmax)$ evaluations, and
$\nmin$ is closed-form.
\end{Proof}

\begin{Corollary}[Infeasibility certificate]
\label{cor:infeasible}
If $\nmin>\nmax$ then no validator-set size satisfies both requirements, and the
deployment must change its architecture --- sharding, delegation, or a weaker
safety target --- rather than its provisioning.
\end{Corollary}

Chance-constrained programming is classical~\cite{charnes1959,nemirovski2007};
what is specific here is that the two constraints are monotone in opposite
directions, so the feasible set is bounded on both sides and infeasibility is
certified rather than merely observed. We are not aware of an equivalent
statement for validator-set sizing. The resulting procedure --- evaluate
$\cc^{\star}$, compute $\nmin$ in closed form, reject immediately if
$h(\nmin)<0$ by Corollary~\ref{cor:infeasible}, otherwise bisect on
$[\nmin,N_{\max}]$ for $\nmax$ and return $\nopt=\nmin$ --- is what
Section~\ref{sec:case} applies.
